\section{A finite tree with nested periods}\label{sec:finite}

We first construct words over a finite alphabet that grows with the
tree. At every internal node, one vector advances the rows cheaply and
a multiple of it has a prescribed smaller period. The latter vector
separates words for every state. We then use Fourier analysis to show
that an $\ellone$ map cannot preserve both kinds of displacement with
small distortion. The short binary code from
Lemma~\ref{lem:delimiter} will then complete the first proof of
Theorem~\ref{thm:binary}.

\subsection{Periods and word construction}

Fix a sufficiently large integer $k$ and put
\begin{equation}\label{eq:finite-parameters}
 m=20k,\qquad T=k^{20k}.
\end{equation}
Use the full ordered $m$-ary tree of depth $k$, with root $o$ at depth
zero and children $v1,\ldots,vm$ at each internal node $v$.
By Lemma~\ref{lem:primes}, choose primes
$p_{j,i}\in[k^4,k^5]$ distinct over all pairs
$0\leq j<k$, $1\leq i\leq m$. Nodes at the same depth use the same
list. Define their periods by
\[
 Q_o=2,\qquad Q_{vi}=Q_vp_{j,i}
 \quad\text{when $v$ has depth $j$}.
\]
All these primes are odd for large $k$.

On the finite abelian group
\[
 A=\prod_{u\text{ leaf}}\Z/Q_u\Z
\]
define vectors supported on their respective subtrees:
\begin{equation}\label{eq:finite-vectors}
 \begin{aligned}
 w_u&=e_u &&(u\text{ leaf}),\\
 u_v&=\sum_{i=1}^m w_{vi},\qquad
 R_v=\prod_{i=1}^m p_{j,i},\qquad
 w_v=R_vu_v
 &&(v\text{ internal at depth }j).
 \end{aligned}
\end{equation}
Here $e_u$ increments the coordinate of leaf $u$.

\begin{lemma}[Orders of the node vectors]\label{lem:orders}
For every node $v$, $\ord(w_v)=Q_v$. At every internal node,
\[
 \ord(u_v)=Q_vR_v,\qquad \gcd(Q_v,R_v)=1.
\]
\end{lemma}
\begin{proof}
At a leaf, $e_u$ has order $Q_u$. Suppose the assertion holds at the
children of an internal node $v$. Their vectors have disjoint supports
and orders $Q_vp_{j,i}$, so their sum has order
\[
 \ord(u_v)=\operatorname{lcm}_{1\leq i\leq m}(Q_vp_{j,i})=Q_vR_v.
\]
The primes dividing $Q_v$ come from earlier depths, together with two,
and none divides $R_v$. Multiplication by $R_v$ therefore leaves a
vector of order $Q_v$, namely $w_v$.
\end{proof}

For $x\in A$, the leaf word is the single symbol $G_u(x)=(u,x_u)$.
At an internal node, define
\begin{equation}\label{eq:finite-word}
 G_v(x)=\prod_{a=0}^{T-1}
 \bigl(G_{v1}(x+a u_v)\cdots G_{vm}(x+a u_v)\bigr),
\end{equation}
where the product concatenates the rows in increasing $a$.
Only leaf coordinates below $v$ affect $G_v$. At depth $j$ and height
$h=k-j$, its length is $L_v=(mT)^h$. Write
\[
 M=(mT)^k,\qquad G=G_o.
\]
Every symbol retains its leaf label. Thus the alphabets occurring
below different children of a node are disjoint, exactly as required by
Lemma~\ref{lem:tagged}.

\subsection{Cheap translations and uniform separation}

We first bound the cost of advancing rows or changing one leaf. These
bounds hold at every state and will later be averaged over $A$.

\begin{lemma}[Inexpensive finite translations]\label{lem:finite-cost}
For every leaf $u$, every $x\in A$ and every integer $t$,
\[
 \ED(G(x),G(x+t e_u))\leq M/m^k.
\]
If $v$ is internal at depth $j$, then for
every $x\in A$ and integer $1\leq s\leq T$,
\begin{equation}\label{eq:finite-cost}
 \ED(G(x),G(x+s u_v))\leq\frac{2s}{T}\frac{M}{m^j}.
\end{equation}
\end{lemma}
\begin{proof}
A leaf occurs $T^k=M/m^k$ times in the root word. Altering that
coordinate can change only those symbols, so substitutions prove the
first bound.

Within one copy of $G_v$, translation by $s u_v$ changes the row
parameters from $0,\ldots,T-1$ to $s,\ldots,T-1+s$. The two lists share
$T-s$ rows. Deleting the first $s$ rows on one side and inserting the
last $s$ on the other costs at most $2sL_v/T$ edits. A node at depth
$j$ occurs $T^j$ times in $G$, with total length $T^jL_v=M/m^j$.
Each occurrence has an input offset contributed by ancestors. That
offset is the same on both sides and commutes with $s u_v$, so the
local script applies in every occurrence. The word is unchanged outside
them. Summing the costs proves \eqref{eq:finite-cost}.
\end{proof}

We now prove separation for every nonzero residue along the
order-$Q_v$ vector $w_v=R_vu_v$, uniformly in the initial state.

\begin{proposition}[Separation in the finite tree]\label{prop:finite-separation}
For every node $v$ of height $h$, every $x\in A$, and every integer
$t$ with $Q_v\nmid t$,
\begin{equation}\label{eq:finite-separation}
 r_{L_v}\bigl(G_v(x),G_v(x+t w_v)\bigr)\geq4^{-h}L_v.
\end{equation}
\end{proposition}
\begin{proof}
At a leaf, the two residue symbols differ. Suppose the proposition
holds at the children of an internal node $v$, and write
$B=L_v/(mT)$ for their common word length.

Compare any source row $a$ in $G_v(x)$ and any target row $a'$ in
$G_v(x+t w_v)$. In child $vi$, the relative input shift is
\[
 (a'-a+tR_v)w_{vi}.
\]
At least $m/2$ children have $Q_{vi}\nmid a'-a+tR_v$. If $a'=a$,
the assertion holds for every child: $Q_v\nmid t$ and
$\gcd(Q_v,R_v)=1$ imply $Q_v\nmid tR_v$. If $a'\ne a$, divisibility
by $Q_{vi}$ would require $p_{j,i}\mid a'-a$, since
$p_{j,i}\mid R_v$. But
\[
 0<|a'-a|<T=k^{20k}
 \quad\text{and}\quad
 (k^4)^{m/2}=k^{40k}>T.
\]
Thus $m/2$ distinct primes of the required sizes cannot all divide
$a'-a$.

By induction, the sum of the deficits over the child blocks in every
pair of rows is at least
\begin{equation}\label{eq:finite-charge}
 d_h=\frac m2\,4^{-(h-1)}B.
\end{equation}
The child blocks have equal length, fixed order and disjoint alphabets.
Also $d_h\leq(m-1)B$. Lemma~\ref{lem:tagged} therefore gives
\[
 r_{L_v}\bigl(G_v(x),G_v(x+t w_v)\bigr)
 \geq\frac{Td_h}{2}
 =4^{-h}mTB=4^{-h}L_v.
\]
The lemma applies to arbitrary matchings, so this proves the induction
for every state $x$.
\end{proof}

We have obtained the two geometric facts needed for the obstruction:
the root shift $w_o$ is far at every state, while the shifts in
Lemma~\ref{lem:finite-cost} are cheap. The remaining argument compares
their average images under an arbitrary $\ellone$ embedding.

\subsection{Fourier detection of sparse activity}

For a finite abelian group $A$, use uniform probability measure and
write $\widehat A$ for its character group. The following identity
represents an averaged $\ellone$ displacement through cuts and Fourier
coefficients. It is the elementary translation principle underlying the comparison with the
cut methods of Khot--Naor and Krauthgamer--Rabani
\cite{KhotNaor2006,KrauthgamerRabani2009}.

\begin{lemma}[Finite-group displacement]\label{lem:finite-fourier}
If $F:A\to\ellone$ is any map on a finite abelian group, then there
are weights $H(\chi)\geq0$ such that, for every $z\in A$,
\begin{equation}\label{eq:finite-fourier}
 \E_x\|F(x+z)-F(x)\|_1
 =\sum_{\chi\in\widehat A}|1-\chi(z)|^2H(\chi).
\end{equation}
The map $F$ need not be injective.
\end{lemma}
\begin{proof}
Apply Lemma~\ref{lem:cuts} to the finite image of $F$ and pull each
cut back to an indicator $U:A\to\{0,1\}$. Since an indicator difference
has absolute value equal to its squared absolute value, Parseval gives
\[
 \E_x|U(x+z)-U(x)|
 =\sum_{\chi\in\widehat A}
 |\widehat U(\chi)|^2|1-\chi(z)|^2,
 \qquad
 \widehat U(\chi)=\E_x U(x)\overline{\chi(x)}.
\]
Multiply by the cut weights and add. The coefficient of each multiplier
is the nonnegative number
$H(\chi)=\sum_C a_C|\widehat U_C(\chi)|^2$.
\end{proof}

For a character $\chi$, call a node $v$ active when $\chi(w_v)\ne1$.
Let $\wt(\chi)$ be the number of active leaves. The period design has
one useful consequence: activity at a node with few active children is
detected by a short average of cheap shifts. Frequencies that escape
all these tests must have large leaf support.

\begin{proposition}[Finite obstruction]\label{prop:finite-lower}
For all sufficiently large $k$, the word image $\{G(x):x\in A\}$,
with ordinary edit distance, has $\ellone$ distortion at least
$C_1^{-1}(k/4)^k$, for an absolute constant $C_1$.
\end{proposition}
\begin{proof}
Take an embedding of the word image with distortion $D$ and rescale it
so that its induced distance $\rho$ on parameter states satisfies
\[
 \frac1D\ED(G(x),G(y))\leq\rho(x,y)\leq\ED(G(x),G(y)).
\]
The map $x\mapsto G(x)$ need not be injective; composing it with the
embedding still gives a map to which Lemma~\ref{lem:finite-fourier}
applies. Write
\[
 V(z)=\E_x\rho(x,x+z)
 =\sum_\chi |1-\chi(z)|^2H(\chi),\qquad H(\chi)\geq0.
\]

First average over the entire cyclic coordinate of one leaf $u$.
The average of $|1-\chi(t e_u)|^2$ is two if $\chi$ is nontrivial on
that coordinate and zero otherwise. Sum the leaf-cost bound in
Lemma~\ref{lem:finite-cost} over all $m^k$ leaves. It follows that
\begin{equation}\label{eq:finite-leaf-mass}
 2\sum_\chi \wt(\chi)H(\chi)\leq M.
\end{equation}

Now suppose an internal node $v$ at depth $j$ is active for $\chi$ but
has fewer than $k$ active children. Since
$\chi(w_v)=\chi(u_v)^{R_v}\ne1$, the value $\chi(u_v)$ is a nontrivial
root of unity. The inactive children contribute one to
$\chi(u_v)=\prod_i\chi(w_{vi})$, so at least one child is active.
Let $q=\ord(\chi(u_v))$. This order divides the least common multiple
of the active child periods, and
\[
 q\leq Q_v\prod_{i:\,vi\text{ active}}p_{j,i}
 \leq2k^{5j+5(k-1)}\leq2k^{10k}.
\]
Put $J=4k^{10k}$. For large $k$, $J\leq T$ and $q\leq J/2$.
The integers $1,\ldots,J$ contain $\lfloor J/q\rfloor$ complete periods
of this root of unity, with total length at least $J/2$. On a complete
nontrivial period the mean of $|1-\chi(su_v)|^2$ is two. Consequently
\[
 \frac1J\sum_{s=1}^J |1-\chi(su_v)|^2\geq1.
\]
This uses complete periods inside the averaging interval and does not
require $q$ to divide $J$.

Because all Fourier weights are nonnegative, the total weight of
characters with this property at the fixed node $v$ is at most
$J^{-1}\sum_{s=1}^J V(su_v)$. Lemma~\ref{lem:finite-cost} bounds that
average by $(2J/T)M/m^j$. There are $m^j$ nodes at depth $j$.
Summing over nodes and all $k$ internal depths bounds the weight of
characters having any active internal node with fewer than $k$ active
children by
\begin{equation}\label{eq:finite-exceptional-mass}
 2k(J/T)M.
\end{equation}
The sets of characters for different nodes may overlap; nonnegativity
makes their sum an upper bound for the union.

Every other character active at the root has at least $k$ active
children at each active internal node. Iterating down the tree gives
$\wt(\chi)\geq k^k$. By \eqref{eq:finite-leaf-mass}, these characters
have total weight at most $M/(2k^k)$. Root-inactive characters contribute
zero to $V(w_o)$, and any root multiplier is at most four. Hence
\begin{equation}\label{eq:finite-root-upper}
 V(w_o)\leq4M\left(\frac{2kJ}{T}+\frac1{2k^k}\right)
 \leq C_1Mk^{-k}.
\end{equation}
Here $J/T=4k^{-10k}$, so an absolute $C_1$ suffices for all large $k$.

Finally $Q_o=2$. Proposition~\ref{prop:finite-separation} with $t=1$
and \eqref{eq:deficit} imply
$\rho(x,x+w_o)\geq4^{-k}M/D$ at every state. Its average, compared with
\eqref{eq:finite-root-upper}, gives
$D\geq C_1^{-1}(k/4)^k$.
\end{proof}

\subsection{Binary witnesses for every sufficiently large cap}

The finite tree alphabet contains one symbol for each leaf and each
residue in its coordinate group. Since $Q_u\leq2k^{5k}$ at every
leaf, its size $K$ satisfies
\[
 K\leq 2m^k k^{5k}=2(20k)^k k^{5k}.
\]
Choose
\[
 t_k=\left\lceil\log_2\bigl(2(20k)^k k^{5k}\bigr)\right\rceil,
 \qquad B_k=2t_k+3=O(k\log k).
\]
There are enough $t_k$-bit labels for all symbols, so
Lemma~\ref{lem:delimiter} applies with width $B_k$.
Its distortion loss is at most $4B_k$. Hence
Proposition~\ref{prop:finite-lower} supplies a finite binary witness
of the single length
\begin{equation}\label{eq:finite-binary-length}
 d_k=B_kM=B_k[20k\cdot k^{20k}]^k
\end{equation}
whose distortion is at least
\[
 D_k:=\frac{(k/4)^k}{4C_1B_k}.
\]
The separated root pair and the lower bound in
Lemma~\ref{lem:delimiter} ensure that the witness contains at least
two words.

The explicit widths $B_k$ and lengths $M$ increase with $k$, so $d_k$
increases to infinity, independently of which admissible primes were
chosen. Moreover,
\begin{align*}
 \log d_k&=20k^2\log k+k\log(20k)+\log B_k\\
 &=20k^2\log k+O(k\log k),\\
 \log D_k&\geq c'k\log k
\end{align*}
for an absolute $c'>0$ and all sufficiently large $k$.
Given a sufficiently large integer $d$, choose $k$ with
$d_k\leq d<d_{k+1}$. Then
$\log d\leq Ck^2\log k$ and $\log\log d\leq C'\log k$ for absolute
constants. Thus
\[
 \log D_k\geq c\sqrt{\log d\,\log\log d}.
\]
The witness of length $d_k\leq d$ proves
Theorem~\ref{thm:binary} for every sufficiently large cap.
